%!TEX root = ../main.tex

\section*{ЗАКЛЮЧЕНИЕ}
\addcontentsline{toc}{section}{ЗАКЛЮЧЕНИЕ}

Настоящая работа старается дать ответы на определенную часть из числа вопросов, возникших в последнее время при работе научного коллектива с моделью развития канала электрического пробоя в твердом диэлектрике, относящейся к типу диффузной границы. Актуальные вопросы можно разделить на три основных направления:
\begin{itemize}
	\item об адекватности математической модели самой по себе, с точки зрения существования и вида решений дифференциальной задачи; в частности, о корректности описания фазовым полем одномерных включений, о выборе определяющих соотношений;
	\item о моделировании задачи вообще; о выборе разностных аппроксимаций, методов решения конечномерной задачи, о ее обусловленности и <<жесткости>> интерполирующих функций;
	\item о корректности и эффективности работы конкретной программы для моделирования задачи в трехмерной постановке.
\end{itemize}
Настоящая работа отчасти касается всех трех направлений исследования.

С точки зрения математической теории, выведены балансовые соотношения \linebreak для внутренней энергии $W$ и свободной энергии $\Psi$; сформулировано несколько возможных типов граничных условий: помимо уже использующегося постоянного потенциала на обкладках, предложены условия постоянного заряда на обкладках и сохранения заряда в системе. Для различных постановок задачи приведены условия устойчивости положений равновесия вида $\phi \equiv C$. Исследовано обобщение модели, частично уже используемое в практике моделирования, допускающее существование решений у ряда модельных краевых задач и, таким образом, позволяющее описывать фазовым полем включения высшей коразмерности.

Для явной схемы разностной задачи в различных постановках предложены два условия устойчивости. При выводе условий произведен определенный анализ используемых определяющих соотношений модели, то есть функций зависимости свойств среды~$\epsilon$, $\sigma$ от фазового поля $\phi$. Также исследованы три различных метода адаптации расчетного шага по времени в моделировании задачи, один из которых~-- на основе условия устойчивости схемы~-- предложен автором. Проведены вычислительные эксперименты в~программе для упрощенного, одномерного, случая задачи.

Наконец, произведено внедрение наиболее удачного (и простого) метода адаптации шага по времени в разрабатываемую научным коллективом программу для моделирования полной, трехмерной, задачи. В программе с адаптацией выполнены два расчета в~модельных постановках задачи, исследованы их результаты. В проделанных тестовых расчетах адаптация сработала удачно, однако на одном из графиков поведения свободной энергии наблюдается аномалия, что, по-видимому, свидетельствует о~не~вполне корректной программной обработке некоторых переменных системы.

Ряд вопросов, касающихся модели, остаются открытыми. Среди них~-- еще не исследованное поведение свободных зарядов в системе (и корректность их программной обработки); анализ различных вариантов определяющих соотношений модели; дальнейшее выявление и решение проблем, связанных с низкой (линейной) размерностью включений, описываемых фазовым полем в задаче развития канала пробоя. По некоторым из них у научного коллектива есть определенные доводы, некоторые только предстоит исследовать в будущем.